% GENERATED FILE. Edit canonical lessons, then run npm run generate:books.
% canonical-source-hash: fnv1a64:11d4f6649b304904
% canonical-lessons: HI-C392-taiyari, HI-C392-azmana, HI-C392-tariqa, HI-C392-anubhav, HI-C392-sujhna

\chapter{Getting Ready}
\label{ch:hi-getting-ready}
\hlchaptermodality{392}

\begin{chapteropening}
\textbf{By the end of this chapter:} \emph{I can say what to try if the first plan fails.}

Complete the last lesson of chapter 392: I can say what to try if the first plan fails.
\end{chapteropening}

\section[\texorpdfstring{\hi{तैयारी} (taiyārī)}{तैयारी (taiyārī)}]{\hi{तैयारी} (taiyārī) — preparation}

\label{lesson:HI-C392-taiyari}



\begin{quote}
Before the new one: say the Hindi for without a break, then the Hindi for to repeat.
\end{quote}

\subsection*{You'll want to know: \hi{तैयारी}}
\textbf{\hi{तैयारी}} — \emph{taiyārī} — \textquotedblleft{}preparation\textquotedblright{}.

\begin{quote}
\hi{अगर} \hi{तैयारी} \hi{अच्छी} \hi{होगी} \hi{तो} \hi{परीक्षा} \hi{से} \hi{डर} \hi{नहीं} \hi{लगेगा।}
\end{quote}

\begin{quote}
\emph{agar taiyārī acchī hogī to parīkṣā se ḍar nahīṃ lagegā.}
\end{quote}

\begin{quote}
If your preparation is good, you won't be afraid of the exam.
\end{quote}

Feminine.

\begin{cousinweb}
\textbf{\hi{तैयारी}} (\emph{taiyārī}) is usually traced to Arabic \emph{tayyār}, \textquotedblleft{}ready\textquotedblright{}, which reached Hindi through Persian. The Hindi noun adds the ending \emph{-ī}.
\end{cousinweb}

\begin{grammarlens}[title={What you've built}]
One more word for saying what happens if.
\end{grammarlens}

\subsection*{Your turn}
\begin{itemize}
\raggedright
  \item \emph{Say it:} \emph{taiyārī}
  \item \emph{Say it:} \emph{taiyārī}, once more
  \item \emph{Say it:} \emph{agar taiyārī acchī hogī to parīkṣā se ḍar nahīṃ lagegā.} — If your preparation is good, you won't be afraid of the exam.
  \item \emph{Recall:} say the Hindi for without a break, then the Hindi for to repeat, then say \emph{taiyārī} again
\end{itemize}

\begin{tcolorbox}[breakable,colback=teal!4,colframe=teal!35!black,title={Before you move on}]
What does \hi{तैयारी} mean? (\textquotedblleft{}Preparation\textquotedblright{}.) Say it once more.
\end{tcolorbox}
\section[\texorpdfstring{\hi{आज़माना} (āzmānā)}{आज़माना (āzmānā)}]{\hi{आज़माना} (āzmānā) — to try out, to test}

\label{lesson:HI-C392-azmana}



\begin{quote}
Before the new one: say the Hindi for to repeat, then the Hindi for preparation.
\end{quote}

\subsection*{You'll want to know: \hi{आज़माना}}
\textbf{\hi{आज़माना}} — \emph{āzmānā} — \textquotedblleft{}to try out, to test\textquotedblright{}.

\begin{quote}
\hi{अगर} \hi{यह} \hi{दवा} \hi{काम} \hi{नहीं} \hi{करेगी} \hi{तो} \hi{डाक्टर} \hi{दूसरी} \hi{दवा} \hi{आज़माएगा।}
\end{quote}

\begin{quote}
\emph{agar yah davā kām nahīṃ karegī to ḍākṭar dūsrī davā āzmāegā.}
\end{quote}

\begin{quote}
If this medicine doesn't work, the doctor will try another one.
\end{quote}

Past, with \hi{ने} (he or she …): \hi{उसने} \hi{नई} \hi{दवा} \hi{आज़माई}, \emph{usne naī davā āzmāī}.

\begin{cousinweb}
\textbf{\hi{आज़माना}} (\emph{āzmānā}) is usually traced to Persian \emph{āzmā-}, \textquotedblleft{}test\textquotedblright{}. This is the stem of \emph{āzmūdan}, \textquotedblleft{}to test\textquotedblright{}.
\end{cousinweb}

\begin{grammarlens}[title={What you've built}]
One more word for saying what happens if.
\end{grammarlens}

\subsection*{Your turn}
\begin{itemize}
\raggedright
  \item \emph{Say it:} \emph{āzmānā}
  \item \emph{Say it:} \emph{āzmānā}, once more
  \item \emph{Say it:} \emph{agar yah davā kām nahīṃ karegī to ḍākṭar dūsrī davā āzmāegā.} — If this medicine doesn't work, the doctor will try another one.
  \item \emph{Recall:} say the Hindi for to repeat, then the Hindi for preparation, then say \emph{āzmānā} again
\end{itemize}

\begin{tcolorbox}[breakable,colback=teal!4,colframe=teal!35!black,title={Before you move on}]
What does \hi{आज़माना} mean? (\textquotedblleft{}To try out, to test\textquotedblright{}.) Say it once more.
\end{tcolorbox}
\section[\texorpdfstring{\hi{तरीक़ा} (tarīqā)}{तरीक़ा (tarīqā)}]{\hi{तरीक़ा} (tarīqā) — a method, a way (of doing)}

\label{lesson:HI-C392-tariqa}



\begin{quote}
Before the new one: say the Hindi for preparation, then the Hindi for to try out.
\end{quote}

\subsection*{You'll want to know: \hi{तरीक़ा}}
\textbf{\hi{तरीक़ा}} — \emph{tarīqā} — \textquotedblleft{}a method, a way (of doing)\textquotedblright{}.

\begin{quote}
\hi{अगर} \hi{यह} \hi{तरीक़ा} \hi{काम} \hi{नहीं} \hi{करेगा} \hi{तो} \hi{हम} \hi{कोई} \hi{दूसरा} \hi{तरीक़ा} \hi{आज़माएँगे।}
\end{quote}

\begin{quote}
\emph{agar yah tarīqā kām nahīṃ karegā to ham koī dūsrā tarīqā āzmāeṅge.}
\end{quote}

\begin{quote}
If this method doesn't work, we'll try another one.
\end{quote}

Masculine.

\begin{cousinweb}
\textbf{\hi{तरीक़ा}} (\emph{tarīqā}) is usually traced to Arabic \emph{ṭarīqa}, \textquotedblleft{}a way, a path\textquotedblright{}, which reached Hindi through Persian.
\end{cousinweb}

\begin{grammarlens}[title={What you've built}]
One more word for saying what happens if.
\end{grammarlens}

\subsection*{Your turn}
\begin{itemize}
\raggedright
  \item \emph{Say it:} \emph{tarīqā}
  \item \emph{Say it:} \emph{tarīqā}, once more
  \item \emph{Say it:} \emph{agar yah tarīqā kām nahīṃ karegā to ham koī dūsrā tarīqā āzmāeṅge.} — If this method doesn't work, we'll try another one.
  \item \emph{Recall:} say the Hindi for preparation, then the Hindi for to try out, then say \emph{tarīqā} again
\end{itemize}

\begin{tcolorbox}[breakable,colback=teal!4,colframe=teal!35!black,title={Before you move on}]
What does \hi{तरीक़ा} mean? (\textquotedblleft{}A method, a way (of doing)\textquotedblright{}.) Say it once more.
\end{tcolorbox}
\section[\texorpdfstring{\hi{अनुभव} (anubhav)}{अनुभव (anubhav)}]{\hi{अनुभव} (anubhav) — experience}

\label{lesson:HI-C392-anubhav}



\begin{quote}
Before the new one: say the Hindi for to try out, then the Hindi for a method.
\end{quote}

\subsection*{You'll want to know: \hi{अनुभव}}
\textbf{\hi{अनुभव}} — \emph{anubhav} — \textquotedblleft{}experience\textquotedblright{}.

\begin{quote}
\hi{अगर} \hi{तुम्हें} \hi{अनुभव} \hi{होगा} \hi{तो} \hi{यह} \hi{नौकरी} \hi{जल्दी} \hi{मिलेगी।}
\end{quote}

\begin{quote}
\emph{agar tumhẽ anubhav hogā to yah naukrī jaldī milegī.}
\end{quote}

\begin{quote}
If you have experience, you'll get this job quickly.
\end{quote}

Masculine.

\begin{cousinweb}
\textbf{\hi{अनुभव}} (\emph{anubhav}) is a tatsama, borrowed from Sanskrit \emph{anubhava}, \textquotedblleft{}perception, experience\textquotedblright{} (Monier-Williams).
\end{cousinweb}

\begin{grammarlens}[title={What you've built}]
One more word for saying what happens if.
\end{grammarlens}

\subsection*{Your turn}
\begin{itemize}
\raggedright
  \item \emph{Say it:} \emph{anubhav}
  \item \emph{Say it:} \emph{anubhav}, once more
  \item \emph{Say it:} \emph{agar tumhẽ anubhav hogā to yah naukrī jaldī milegī.} — If you have experience, you'll get this job quickly.
  \item \emph{Recall:} say the Hindi for to try out, then the Hindi for a method, then say \emph{anubhav} again
\end{itemize}

\begin{tcolorbox}[breakable,colback=teal!4,colframe=teal!35!black,title={Before you move on}]
What does \hi{अनुभव} mean? (\textquotedblleft{}Experience\textquotedblright{}.) Say it once more.
\end{tcolorbox}
\section[\texorpdfstring{\hi{सूझना} (sūjhnā)}{सूझना (sūjhnā)}]{\hi{सूझना} (sūjhnā) — to occur to (someone), to strike}

\label{lesson:HI-C392-sujhna}



\begin{quote}
Before the new one: say the Hindi for a method, then the Hindi for experience.
\end{quote}

\subsection*{You'll want to know: \hi{सूझना}}
\textbf{\hi{सूझना}} — \emph{sūjhnā} — \textquotedblleft{}to occur to (someone), to strike\textquotedblright{}.

\begin{quote}
\hi{अगर} \hi{तुम्हें} \hi{कोई} \hi{उपाय} \hi{सूझेगा} \hi{तो} \hi{मुझे} \hi{बताओ।}
\end{quote}

\begin{quote}
\emph{agar tumhẽ koī upāy sūjhegā to mujhe batāo.}
\end{quote}

\begin{quote}
If a way out occurs to you, tell me.
\end{quote}

Past (it …): \hi{वह} \hi{उपाय} \hi{उसे} \hi{सूझा}, \emph{vah upāy use sūjhā}.

\begin{cousinweb}
\textbf{\hi{सूझना}} (\emph{sūjhnā}) is a tadbhava, inherited from Sanskrit \emph{śudhyati}, \textquotedblleft{}becomes clear\textquotedblright{} (CDIAL 12527).
\end{cousinweb}

\begin{grammarlens}[title={What you've built}]
One more word for saying what happens if.
\end{grammarlens}

\subsection*{Your turn}
\begin{itemize}
\raggedright
  \item \emph{Say it:} \emph{sūjhnā}
  \item \emph{Say it:} \emph{sūjhnā}, once more
  \item \emph{Say it:} \emph{agar tumhẽ koī upāy sūjhegā to mujhe batāo.} — If a way out occurs to you, tell me.
  \item \emph{Recall:} say the Hindi for a method, then the Hindi for experience, then say \emph{sūjhnā} again
\end{itemize}

\begin{tcolorbox}[breakable,colback=teal!4,colframe=teal!35!black,title={Before you move on}]
What does \hi{सूझना} mean? (\textquotedblleft{}To occur to (someone), to strike\textquotedblright{}.) Say it once more.
\end{tcolorbox}
